\section{Alignment Function}
\begin{frame}{Alignment Function}
  \begin{itemize}
    \item The alignment function is the component of attention that determines how much focus to place on different parts of the input sequence.
    \item It computes a score for each input token relative to the current output token being generated.
    \item The scores are then normalized (often using softmax) to create a probability distribution, indicating the relative importance of each input token.
  \end{itemize}
\end{frame}

\begin{frame}{Importance of Alignment Function}
  \justifying
  The choice of alignment function can significantly impact the model’s performance, especially in tasks with long sequences or complex dependencies.

  \vspace{1em}
  \small
  \textbf{Example:} In a translation task, the alignment function helps the model determine which words in the source language should be emphasized when generating each word in the target language. For instance, when translating “The cat sat on the mat,” the model might focus more on “cat” when generating “gato” in Spanish, while still considering “sat” and “on” to maintain the context.
\end{frame}

\begin{frame}{Types of Alignment Functions}
  Here $\mathbf{q}$ is the query formed from the decoder state and $\mathbf{k}$ the key formed from an encoder state; both are made explicit in the next section.
  \begin{itemize}
    \item[\textcolor{red}{$\blacktriangleright$}] \textbf{Dot Product}:
      \begin{equation*}
        \text{score} = \mathbf{q}^\top \mathbf{k}
      \end{equation*}
    \item[\textcolor{red}{$\blacktriangleright$}] \textbf{Scaled Dot Product}:
      \begin{equation*}
        \text{score} = \frac{\mathbf{q}^\top \mathbf{k}}{\sqrt{d_k}}
      \end{equation*}
    \item[\textcolor{red}{$\blacktriangleright$}] \textbf{Additive (Bahdanau) Attention}:
      \begin{equation*}
        \text{score} = \mathbf{v}^\top \tanh(W_1 h_i + W_2 s_{t-1})
      \end{equation*}
    \item[\textcolor{red}{$\blacktriangleright$}] \textbf{General}:
      \begin{equation*}
        \text{score} = \mathbf{q}^\top W_a \mathbf{k}
      \end{equation*}
  \end{itemize}
\end{frame}

\begin{frame}{Choosing an Alignment Function}
  \begin{itemize}
    \item[\textcolor{red}{$\blacktriangleright$}] \textbf{Dot product} adds no parameters and reduces to a single matrix product, so it is the cheapest, but it requires $\mathbf{q}$ and $\mathbf{k}$ to have the same dimension.
    \item[\textcolor{red}{$\blacktriangleright$}] \textbf{Scaling by $\sqrt{d_k}$} matters because, for independent components with zero mean and unit variance, $\mathbf{q}^\top\mathbf{k}$ has variance $d_k$. Left unscaled, the scores grow with dimension and push the softmax into a region where its gradients are vanishingly small.
    \item[\textcolor{red}{$\blacktriangleright$}] \textbf{Additive attention} passes the pair through a hidden layer, so query and key need not share a dimension, at the cost of an extra weight matrix and a nonlinearity per score.
    \item[\textcolor{red}{$\blacktriangleright$}] \textbf{General} sits between the two: one learned matrix $W_a$ bridges mismatched dimensions while keeping the score a single product.
  \end{itemize}
\end{frame}
